
\section{Open Research Questions for Human-like Interaction}
\label{sec:prospects}

The central aim of speech interaction is to build intelligence for communicating with people. This intelligence is expressed through both conversational content and rhythm: understanding and formulating what to say, and organizing how an exchange unfolds. Our research agenda therefore concerns the formation of interaction capabilities and the human experience they enable, extending from context-dependent dialogue to learning through speech, participation in real environments, and collaboration with execution systems. The following questions articulate this direction and the mechanisms that remain to be understood.

\textbf{How can models learn the context-conditioned organization of natural interaction?} Interruption, backchanneling, turn transitions, and proactive questions are manifestations of a broader ability to organize an exchange. Studying them as isolated skills leaves open how they work together to make interaction natural and comfortable. A model must integrate acoustic and semantic evidence with conversational history, environmental conditions, task context, and user preferences. It must determine when to enter, hold, or yield the floor, how much to say, and when a question or suggestion would help the user. The same pause or utterance can call for different behavior in different situations. Research should examine how these conditions jointly shape conversational decisions and how local decisions contribute to the experience of an entire exchange. Learning from contextualized human dialogue and human–system interaction can connect behavior patterns with their communicative effects. Evaluation should combine controlled changes in context with complete interactions, examining ease, comfort, continuity, and the feeling of being understood. Skill-specific measurements can diagnose mechanisms, while the organizing objective remains coherent, adaptive interaction.

\textbf{How can learning from speech give rise to language and interaction intelligence?} Language is not confined to writing. Speech carries linguistic structure through temporally unfolding expression and places language within communication, shared situations, and responses from others. It should therefore be investigated as a source of learning language and intelligence, beyond an input representation to be aligned with a pretrained semantic space. A fundamental question is how lexical, grammatical, semantic, and pragmatic knowledge can be learned through speech experience and combined into generalizable interaction competence. This includes understanding intentions, reasoning about context, and coordinating conversational content with rhythm. A nearer-term direction asks how speech experience can enhance the intelligence of existing language models, rather than only add speech recognition and generation interfaces. These directions motivate learning objectives and experience structures that expose models to language in use and to the consequences of their contributions. Evidence should examine new linguistic and interactive situations, including whether the resulting abilities extend beyond the training scenarios. Initialization, auxiliary text, representations, and feedback provide experimental choices for investigating this process. The central target is the contribution of speech experience to intelligence; the role of each learning ingredient should be established through that target.

\textbf{How can a single-user assistant become a situated participant in dialogue?} Speech models intended for real environments must understand their own place within an exchange. In homes, meetings, and public settings, hearing intelligible speech does not automatically mean being addressed. Near- and far-field voices, simultaneous discussion, side conversations, and changing roles require a model to infer participation relationships and conversational rights. Perception provides evidence about who spoke, where speech originated, and what was said; interaction intelligence must determine whether to respond, remain quiet, join a discussion, or contribute proactively. A dialogue participant also needs continuity in its role and an understanding of the topic and people it is engaging with. Research should connect auditory scene understanding with addressee inference, floor management, topic organization, and individual and shared context. Realistic tests should include invitations to participate, competing requests, people joining or leaving, and opportunities for helpful contributions. The desired capability is a situated form of agency through which the model can be available and constructive without treating every nearby utterance as its own instruction. Its quality should be assessed through both appropriate participation and the experience of the people involved.

\textbf{How should interaction and execution intelligence collaborate to serve people?} We take the quality of the human experience as the organizing objective of the interaction model. Its responsibility is to understand users, help them express and refine their needs, guide the exchange, and make services accessible through appropriate content and rhythm. Execution models, tools, and agents primarily undertake the reasoning and actions required to carry out the resulting work. This division gives interaction research its own objective: making people feel understood and effectively supported throughout the service process. The frontend consequently requires substantive understanding and reasoning about people and communication, while the backend provides capabilities for solving tasks and reporting what occurred. Research should determine which abilities and responsibilities belong to each role and how goals, constraints, uncertainty, progress, and results should pass between them. These roles can be implemented through shared models, separate models, or collaborations with coordinating runtimes. Their organization should preserve the frontend’s responsibility for human experience. Honest explanations of capability, progress, failure, and outcomes are part of that experience. Evaluation should examine whether the arrangement reduces communicative effort, supports goal formation, and provides natural, comfortable guidance, with execution competence and faithful reporting supporting those judgments. Effective coordination would allow advances in task capability to become services that people can readily understand and use.

Together, these questions define a program in which speech is a medium for learning intelligence, conversation is organized around people, and interaction models acquire a meaningful role in real settings. Progress depends on linking these capabilities to the experiences they create and the mechanisms through which they are learned and coordinated.



